% =====================================================================
%  Wstęp pracy inżynierskiej — dwie RZECZOWE propozycje (wersja 2, 28.10.2026)
%  Styl: sprawozdawczy, bez pytań retorycznych i ozdobników.
%  \cite{TODO-...} = miejsca, gdzie trzeba podpiąć źródło z bibliografii.
%  Poprzednie propozycje: wstep_zagajenie_opcje.tex (bez zmian).
% =====================================================================

% =====================================================================
%  OPCJA C — od analizy materiałów do warunków misji kosmicznych
% =====================================================================
\chapter{Wstęp}

Analiza składu pierwiastkowego jest jednym z~podstawowych badań
materiałowych. W~warunkach laboratoryjnych wykonuje się ją najczęściej
metodą spektroskopii EDS (ang.~\textit{Energy Dispersive Spectroscopy})
w~elektronowym mikroskopie skaningowym. Wiązka elektronów pada
na~próbkę, a~zawarte w~niej pierwiastki emitują promieniowanie X
o~energiach charakterystycznych dla każdego z~nich. Spektrometr mierzy
energie fotonów, co pozwala ustalić skład próbki.

Ten sam rodzaj badań jest potrzebny w~misjach kosmicznych, w~których
określa się skład skał i~gruntu na~powierzchni Księżyca, Marsa czy
planetoid. Sprowadzanie próbek na~Ziemię jest kosztowne i~możliwe tylko
w~nielicznych misjach, dlatego analizę wykonuje się na~miejscu,
przyrządami zamontowanymi na~lądowniku lub łaziku~\cite{TODO-insitu}.
Stosowane obecnie przyrządy, takie jak spektrometry APXS czy PIXL,
wzbudzają promieniowanie X~próbki źródłem radioizotopowym lub lampą
rentgenowską~\cite{TODO-apxs,TODO-pixl}.

Wzbudzanie próbki bezpośrednio wiązką elektronów, tak jak w~EDS,
nie jest w~tych zastosowaniach stosowane. Wymaga ono działa
elektronowego, które w~mikroskopie wraz z~kolumną elektronooptyczną
i~układem próżniowym ma znaczne wymiary, masę i~pobór mocy.
W~przyrządach kosmicznych są to parametry ograniczone budżetem misji.
Na~powierzchni ciał pozbawionych atmosfery panuje natomiast naturalna
próżnia, w~której wiązka elektronów może propagować bez osobnej komory.

Zmniejszenie wymiarów źródła elektronów umożliwia technologia MEMS
(ang.~\textit{Micro-Electro-Mechanical Systems}), w~której wytwarza się
struktury o~wymiarach od~mikrometrów do~milimetrów. Na~Politechnice
Wrocławskiej opracowano w~tej technologii miniaturowe źródło
promieniowania rentgenowskiego z~zimną katodą z~nanorurek węglowych
(CNT)~\cite{TODO-urbanski}. Katoda ta nie wymaga elementu grzejnego,
a~cała struktura krzemowo-szklana ma wymiary rzędu centymetrów i~masę
rzędu gramów. W~niniejszej pracy rozwiązanie to wykorzystano do~budowy
miniaturowego analizatora EDS, w~którym wiązka elektronów,
ogniskowana polem magnetycznym, pada bezpośrednio na~badaną próbkę.

Tematyka pracy wpisuje się w~dyscyplinę automatyka, elektronika,
elektrotechnika i~technologie kosmiczne, ponieważ łączy elektronikę
układów zasilania źródła elektronów i~akwizycji widm, elektrotechnikę
wysokich napięć przyspieszających i~ogniskowania wiązki oraz
projektowanie przyrządu przeznaczonego do~misji kosmicznych.

Praca jest również zgodna z~profilem kierunku Inteligentna Elektronika.
Analizator EDS można traktować jako czujnik składu pierwiastkowego.
Wielkością mierzoną jest skład chemiczny próbki, a~sygnałem wyjściowym
widmo energii fotonów promieniowania X. Podobnie jak w~innych układach
czujnikowych, można w~nim wyróżnić element wzbudzający, czyli źródło
elektronów, element przetwarzający, czyli detektor z~torem akwizycji,
oraz układ zasilania i~sterowania, który ustala warunki pracy źródła,
m.in. napięcie ekstrakcyjne i~przyspieszające.

W~docelowym przyrządzie kosmicznym taki czujnik pracowałby samodzielnie,
bez bieżącego udziału operatora, ze względu na~opóźnienia łączności
z~Ziemią. Wymaga to elektroniki, która steruje emisją elektronów,
kontroluje stabilność prądu wiązki, rejestruje widmo i~wstępnie je
przetwarza, np.~rozpoznając linie charakterystyczne pierwiastków przed
wysłaniem danych. Zagadnienia te łączą mikrosystemy MEMS z~układami
pomiarowymi i~przetwarzaniem sygnałów, co odpowiada zakresowi kierunku
Inteligentna Elektronika.


% =====================================================================
%  OPCJA D — od wymagań stawianych przyrządom kosmicznym
% =====================================================================
\chapter{Wstęp}

Przyrządy przeznaczone do~misji kosmicznych muszą spełniać wymagania,
które w~aparaturze laboratoryjnej mają drugorzędne znaczenie. Każdy
gram masy i~każdy wat mocy zwiększają koszt misji, a~miejsce
na~pokładzie lądownika lub łazika jest dzielone między wiele
instrumentów~\cite{TODO-budzet}. Z~tego powodu w~technice kosmicznej
dąży się do~miniaturyzacji przyrządów przy zachowaniu ich funkcji
pomiarowych.

Jednym z~podstawowych zadań misji planetarnych jest określenie składu
pierwiastkowego skał i~gruntu. Pozwala to wnioskować o~pochodzeniu
i~historii badanego ciała oraz ocenić dostępność surowców, np.~na
Księżycu~\cite{TODO-sklad}. W~laboratoriach skład materiałów bada się
powszechnie metodą spektroskopii EDS
(ang.~\textit{Energy Dispersive Spectroscopy}). Wiązka elektronów
pada na~próbkę, a~zawarte w~niej pierwiastki emitują promieniowanie X
o~energiach charakterystycznych dla każdego z~nich, które mierzy
spektrometr.

Przeniesienie tej metody na~powierzchnię innego ciała niebieskiego
ograniczają przede wszystkim wymiary i~pobór mocy działa elektronowego.
Klasyczne działo z~żarzoną katodą i~kolumną elektronooptyczną jest
elementem mikroskopu skaningowego i~nie spełnia wymagań przyrządu
kosmicznego. Jednocześnie na~ciałach pozbawionych atmosfery, takich jak
Księżyc czy planetoidy, naturalna próżnia umożliwia propagację wiązki
elektronów bez osobnej komory, co czyni EDS metodą atrakcyjną dla
takich misji.

Ograniczenie wymiarów i~poboru mocy źródła elektronów umożliwia
technologia MEMS (ang.~\textit{Micro-Electro-Mechanical Systems}).
Na~Politechnice Wrocławskiej opracowano w~tej technologii miniaturowe
źródło promieniowania rentgenowskiego z~zimną katodą z~nanorurek
węglowych (CNT)~\cite{TODO-urbanski}. Zimna emisja polowa nie wymaga
podgrzewania katody, a~struktura krzemowo-szklana ma wymiary rzędu
centymetrów i~masę rzędu gramów. Niniejsza praca dotyczy wykorzystania
tego rozwiązania do~budowy miniaturowego analizatora EDS, w~którym
wiązka elektronów, ogniskowana polem magnetycznym, pada bezpośrednio
na~badaną próbkę.

Tematyka pracy wpisuje się w~dyscyplinę automatyka, elektronika,
elektrotechnika i~technologie kosmiczne, ponieważ łączy elektronikę
układów zasilania źródła elektronów i~akwizycji widm, elektrotechnikę
wysokich napięć przyspieszających i~ogniskowania wiązki oraz
projektowanie przyrządu przeznaczonego do~misji kosmicznych.
